%!TEX root = ./springer-mono-format/book.tex
%!TEX root = ./thesis-trier.tex


\begin{figure}[htb]
	\centering
	
	%--------------------------------------------------------------- row 1
	\begin{subfigure}[t]{.48\textwidth}
		\centering
		\scalebox{0.8}{%
			\begin{tikzpicture}[
				baseline=(current bounding box.north),
				node distance=2.4em
				]
				\node[state]              (1) {};
				\node[state, right=of 1]  (2) {};
				\node[state, right=of 2]  (3) {};
				
				\path
				(1) edge node[above] {$a$} (2)
				(2) edge node[above] {$b$} (3);
			\end{tikzpicture}%
		}
		\caption{Three states, with a trap state.}
		\label{prelim:fig:three-state-sync-binary-SI-a}
	\end{subfigure}\hfill
	\begin{subfigure}[t]{.48\textwidth}
		\centering
		\scalebox{0.8}{%
			\begin{tikzpicture}[
				baseline=(current bounding box.north),
				node distance=2.4em
				]
				\node[state]              (1) {};
				\node[state, left=of 1]   (2) {};
				\node[state, right=of 1]  (3) {};
				
				\path
				(2) edge node[above] {$a$} (1)
				(3) edge node[above] {$b$} (1);
			\end{tikzpicture}%
		}
		\caption{Three states, both letters merging.}
		\label{prelim:fig:three-state-sync-binary-SI-b}
	\end{subfigure}
	
	\medskip
	
	%--------------------------------------------------------------- row 2
	\begin{subfigure}[t]{.48\textwidth}
		\centering
		\scalebox{0.8}{%
			\begin{tikzpicture}[
				%baseline=(current bounding box.north),
				node distance=4.5em
				]
				\def\R{7em}
				\useasboundingbox (-7.8em,-8.0em) rectangle (11.6em,8.0em);
				
				\foreach \name/\angle/\text in
				{A1/0/s,A2/45/,A3/90/,A4/135/,A5/180/,A6/225/,A8/315/}
				{
					\node[state] (\name) at (\angle:\R) {$\text$};
				}
				
				\node[state,right of=A1] (t) {$t$};
				
				
				
				\path
				(A6) edge node[auto] {$b$} (A5)
				(A5) edge node[auto] {$b$} (A4)
				(A4) edge node[auto] {$b$} (A3)
				(A3) edge node[auto] {$b$} (A2)
				(A2) edge node[auto] {$b$} (A1)
				(A1) edge node[auto] {$b$} (A8)
				(A1) edge node[above] {$a$} (t);
				%(t)  edge[loop above] node {$a,b$} (t);
				
				\node at (270:\R) {\large$\cdots$};
			\end{tikzpicture}%
		}
		\caption{Binary synchronizing SI-automaton with a trap state.}
		\label{prelim:fig:binarySI-sync-trap}
	\end{subfigure}\hfill
	\begin{subfigure}[t]{.48\textwidth}
		\centering
		\scalebox{0.8}{%
			\begin{tikzpicture}[
				%baseline=(current bounding box.north),
				]
				\useasboundingbox (-7.8em,-8.0em) rectangle (11.6em,8.0em);
				\def\R{7em}
				
				\foreach \name/\angle/\text in
				{A1/0/,A2/45/,A3/90/,A4/135/,A5/180/t,A6/225/,A7/270/,A8/315/s}
				{
					\node[state] (\name) at (\angle:\R) {$\text$};
				}			
				
				\path
				(A1) edge node[auto] {$b$} (A8)
				(A7) edge node[auto] {$b$} (A6)
				(A6) edge node[auto] {$b$} (A5)
				(A5) edge node[auto] {$b$} (A4)
				(A3) edge node[auto] {$b$} (A2)
				(A2) edge node[auto] {$b$} (A1);
				
                \draw (A8) edge [auto] node {$a$} (A5);
				
				\node[rotate=22.5] at (292.5:\R) {\large$\cdots$};
				\node[rotate=202.5] at (112.5:\R) {\large$\cdots$};
			\end{tikzpicture}%
		}
		\caption{Strongly connected binary synchronizing SI-automaton where
			$\applytrans{\delta}{s}{b^d}=t$.}
		\label{prelim:fig:binarySI-sync-SCC}
	\end{subfigure}
	
	\caption[Synchronizing binary SI-automata]{Synchronizing binary SI-automata (self-loops omitted)
		as listed in \cref{prop:structure_simp_idem_binary}.
		Note that, in \cref{prelim:fig:binarySI-sync-SCC}, for $d=1$ we obtain the \Cerny family.}
	\label{prelim:fig:three-state-sync-binary-SI}
	\label{fig:theorem}
\end{figure}